\documentclass[11pt,a4paper]{article}

\usepackage[T1]{fontenc}
\usepackage[utf8]{inputenc}
\usepackage[margin=2.3cm]{geometry}
\usepackage{amsmath,amssymb}
\usepackage{microtype}
\usepackage{enumitem}
\usepackage{xcolor}
\usepackage[hidelinks,colorlinks=true,linkcolor=blue!50!black,citecolor=blue!50!black,urlcolor=blue!50!black]{hyperref}
\usepackage{graphicx}
\usepackage{booktabs}

\newcommand{\OBQF}{\textsc{1BQF}}
\newcommand{\bra}[1]{\langle #1 |}
\newcommand{\ket}[1]{| #1 \rangle}
\newcommand{\braket}[2]{\langle #1 | #2 \rangle}

\title{Generalisability of the 1\,bit Quantum Filter (\OBQF{}) Beyond Track Reconstruction\\[2pt]
\large A response to Reviewer~2's ``too tailored to this Hamiltonian'' concern}
\author{LHCb Quantum Tracking working note}
\date{\today}

\begin{document}
\maketitle

\begin{abstract}
\noindent
The reviewer asks whether the spectral-filter trick at the heart of our quantum
track-reconstruction pipeline -- which we will refer to as the
\emph{one-bit quantum filter} (\OBQF{}) -- is specific to the LHCb VELO
Hopfield Hamiltonian, or whether it represents a more general computational
primitive.  This note argues, with citations, that \OBQF{} is a
\emph{problem-class-level} primitive: it is the minimal quantum
linear-system (HHL-type) circuit whenever (i)~the system matrix admits a
\emph{bimodal} spectrum with one cluster of physically irrelevant
``noise'' eigenvalues, (ii)~the right-hand side preparation is uniform or
otherwise gate-cheap, and (iii)~the quantity of interest is a low-weight
expectation value of the solution, not the full solution vector.  We
identify families of problems where these three conditions hold by
construction (associative memory recall, graph-Laplacian / Hopfield-QUBO
optimisation, kernel-method classifiers, fitting with quadratic
regularisers, certain PDE inverse problems, image template-matching), and
delimit where they do not (dense linear regression with arbitrary right
hand sides, problems requiring full state tomography, ill-conditioned
systems with continuous spectra).
\end{abstract}

\section{What the 1\,bit Quantum Filter actually is}
\label{sec:definition}

In the published track-reconstruction work the quantum solver receives a
Hopfield-Hamiltonian linear system $\mathbf{A}\,\mathbf{x} = \mathbf{b}$
constructed from the segment graph of a VELO event, with $\mathbf{A}$ a
real symmetric positive-definite matrix of dimension $N$ equal to the
number of admissible doublet/segment candidates.  The key spectral
observation -- already implicit in the original Hopfield/Stimpfl-Abele
formulation \cite{Hopfield:1982,StimpflAbele:1991,Denby:1988} -- is that
$\mathbf{A}$ has a \emph{bimodal} spectrum:
\begin{itemize}[leftmargin=1.2em,topsep=2pt,itemsep=2pt]
  \item a small cluster of ``signal'' eigenvalues
        $\{\lambda_s\}$, whose eigenvectors are supported on the segments
        belonging to true tracks;
  \item a large cluster of ``noise'' eigenvalues, concentrated around
        a single, analytically known value $\lambda_n \approx \gamma$
        (the conflict-penalty coupling), whose eigenvectors are the
        Hopfield ``basins'' of incompatible-segment configurations.
\end{itemize}

\OBQF{} is the HHL pipeline of \cite{Harrow:2009} reduced to a
\emph{single} clock (phase-estimation) qubit, $n_t = 1$.  Because the
phase register has only two computational-basis states, the controlled
inverse rotation can resolve only \emph{two} eigenvalue bins.  This is
exactly the resolution required: one bin captures the noise cluster and
its amplitude is rejected via the ancilla post-selection, the other bin
captures the signal cluster.  Schematically, with
$\Pi_\mathrm{sig} = \sum_{\lambda \in \{\lambda_s\}} \ket{\lambda}\bra{\lambda}$
and $\Pi_\mathrm{noise} = \mathbb{1} - \Pi_\mathrm{sig}$,
\begin{equation}
  \ket{\psi_{\mathrm{post}}}
  \;\propto\;
  \Pi_\mathrm{sig}\,\ket{\mathbf{b}}
  \;+\;\mathcal{O}\!\bigl(\sqrt{p_\mathrm{leak}}\bigr)\,\Pi_\mathrm{noise}\,\ket{\mathbf{b}},
  \label{eq:filter}
\end{equation}
where $p_\mathrm{leak}$ is set by the spectral gap
$\Delta = |\lambda_n - \lambda_s|$ and the QPE resolution.  The
final answer is read out as the expectation value of an
$\mathcal{O}(1)$-weight observable (in our case the
which-segment-is-active projector on a chosen qubit), not as a full
amplitude vector.

The three structural ingredients that make this work are therefore:
\begin{enumerate}[label=(\textbf{S\arabic*}),leftmargin=2.4em,topsep=2pt,itemsep=2pt]
  \item \textbf{Bimodal spectrum} with one analytically known cluster
        eigenvalue $\lambda_n$ that is irrelevant to the answer.
  \item \textbf{Gate-cheap right-hand side}: $\ket{\mathbf{b}}$
        preparable without QRAM (in our case $\ket{\mathbf{b}}=H^{\otimes
        n_s}\ket{0}$, the uniform superposition).
  \item \textbf{Low-weight observable}: the answer is
        $\bra{\psi}O\ket{\psi}$ with $O$ supported on $\mathcal{O}(1)$
        qubits, not the amplitudes of $\ket{\psi}$ themselves.
\end{enumerate}
None of (S1)--(S3) refers to tracking, to LHCb, or even to physics. Any
problem that admits this triple is an \OBQF{} problem.

\section{Why \OBQF{} is structurally transferable}
\label{sec:why-transferable}

\OBQF{} is best understood as a special case of the broader family of
\emph{quantum eigenvalue / singular-value transformations} (QSVT)
\cite{Gilyen:2019,Martyn:2021}, in which a polynomial $P(\mathbf{A})$ of
the system matrix is implemented coherently, and matrix inversion is
recovered by choosing $P(\lambda) \approx 1/\lambda$.  The full QSVT
construction is overkill if one only needs the projector
$P(\lambda) \approx \mathbf{1}_{\lambda \in [\lambda_s^-,\lambda_s^+]}$
that separates two clusters: a single bit of phase resolution is
sufficient by Kitaev's iterative phase-estimation argument
\cite{Kitaev:1995,Higgins:2007}.  In other words \OBQF{} is the
\emph{cheapest non-trivial QSVT} that still solves a useful problem; it
is interesting precisely because so many useful problems satisfy
(S1)--(S3) and therefore do not require the full polynomial machinery.

The transferability statement is then:
\begin{quote}\itshape
  Whenever a linear-algebraic task can be re-cast so that its system
  matrix has a bimodal spectrum, its right-hand side is gate-cheap, and
  its answer is a low-weight observable, the same \OBQF{} circuit -- up
  to a problem-specific Hamiltonian-simulation block for $e^{-i\mathbf
  A t}$ -- solves it with the same $\mathcal{O}(\sqrt{N}\log N)$ gate
  budget reported in our paper.
\end{quote}
This argument is structural: it reduces the question
``can your method generalise?'' to the
question ``does my problem have a bimodal spectrum?''.  The latter is a
spectral-graph-theory question that can be answered before any quantum
circuit is written, by classical sparse-eigenvalue analysis (in our work
this is done by ARPACK \cite{Lehoucq:1998:arpack}).

\section{Families of problems that satisfy (S1)--(S3)}
\label{sec:families}

We list, in decreasing order of structural similarity to our
tracking application, the problem families to which \OBQF{}
transfers.  In each case we name the source matrix, the
``noise'' eigenvalue, the natural observable, and a representative
reference.

\subsection*{F1. Hopfield-style associative memory and pattern recognition}
The original Hopfield network \cite{Hopfield:1982} stores patterns as
fixed points of a quadratic energy
$E(\mathbf{x}) = -\tfrac12 \mathbf{x}^\top \mathbf{W}\mathbf{x}$, recall is
linear inversion in the basin of a stored pattern, and the
spurious-pattern spectrum is concentrated on a single
loading-fraction-dependent eigenvalue \cite{Amit:1985}.  Quantum
associative-memory proposals \cite{Ventura:2000,Trugenberger:2001} are
exactly in this regime; \OBQF{} provides the missing
inversion/recall step with single-clock-qubit cost.  Track
reconstruction is one instance of this family: any combinatorial
pattern-completion problem with a Hopfield encoder qualifies (e.g.\
content-addressable memory, image inpainting on small
quantum-tractable patches).

\subsection*{F2. Graph-Laplacian / spectral-clustering linear systems}
For a graph with combinatorial Laplacian $\mathbf{L}$, problems such as
seeded community detection, electrical-network effective-resistance
queries, and PageRank-with-restart are linear systems
$(\mathbf{L} + \alpha\mathbf{1})\mathbf{x} = \mathbf{b}$.  The Laplacian
of a graph with $k$ well-separated communities has $k$ small
eigenvalues and a clearly separated bulk \cite{vonLuxburg:2007}; the
bulk eigenvalue acts as $\lambda_n$ and the answer (community label of
a query node, effective resistance, etc.) is a low-weight observable.
QSVT-based approaches to such problems already exist
\cite{Apers:2022,Childs:2017} and \OBQF{} is the natural
single-clock-qubit specialisation.

\subsection*{F3. Quadratic unconstrained binary optimisation (QUBO)
with structured penalties}
Our Hamiltonian is a QUBO with a large conflict penalty
$\gamma$ that pins the
``noise'' eigenvalue.  This is the standard pattern in
QUBO/Ising encodings of constrained combinatorial problems
\cite{Lucas:2014}: max-cut, graph colouring, set cover, vertex cover,
and maximum independent set all have penalty matrices whose spectrum
clusters around the penalty strength.  When the question of interest is
``is variable $i$ active in the dominant low-energy mode?'' (a
low-weight observable on a single qubit), \OBQF{} returns the answer
directly without sampling the full ground state.  This makes \OBQF{}
complementary to QAOA \cite{Farhi:2014} and quantum annealing
\cite{Kadowaki:1998}: instead of preparing the ground state, it
projects onto the low-energy subspace and reads a marginal.

\subsection*{F4. Kernel methods and least-squares classifiers
(LS-SVM, ridge regression, Gaussian-process regression)}
A regularised kernel system $(\mathbf{K} + \lambda\mathbf{1})\boldsymbol\alpha
= \mathbf{y}$ has a spectrum dominated by a few large kernel
eigenvalues plus a long tail near the regulariser $\lambda$
\cite{Williams:2006}; the regulariser pins $\lambda_n$.  When the test
prediction is needed only for a few new query points -- a low-weight
observable on the output register -- this is exactly an \OBQF{} problem.
Quantum kernel/SVM approaches based on full HHL
\cite{Rebentrost:2014:qsvm} can be reduced to single-clock-qubit cost
whenever the kernel is well-conditioned and the prediction observable is
local, which is the setting most often encountered in HEP fast-inference
applications.

\subsection*{F5. Linear PDEs with a known spectral gap}
Discretised elliptic PDEs $-\nabla\cdot(c\nabla u) = f$ on a
quasi-uniform mesh produce a stiffness matrix whose eigenvalues separate
into a small set tied to global modes and a bulk tied to the local
discretisation scale; this gap is the basis of multigrid preconditioning.
Quantum PDE solvers \cite{Berry:2014:qode,Childs:2017,Berry:2017} that
read out only a local functional of $u$ (a flux, an integrated source
term) are \OBQF{}-compatible.  Track extrapolation through a
piecewise-uniform field map -- a directly LHCb-relevant
PDE/ODE-inverse-problem -- is in this class.

\subsection*{F6. Template matching and constrained image segmentation}
When the answer asked of an image is binary -- ``does pattern $T$
appear at position $\mathbf p$?'' -- the underlying linear problem is
again a quadratic projection onto a known template subspace whose
spectrum is bimodal (template / non-template).  This generalises the
``did this pattern of hits form a track?'' question of our paper.
Quantum template-matching circuits already use spectral-projection
ideas \cite{Cai:2015:templatematch} and admit a one-bit reduction.

\bigskip\noindent
The list is not exhaustive but is sufficient to answer Reviewer~2's
question: the spectral-filter trick is not Hamiltonian-specific; it
works whenever the spectrum has the right shape and the quantity of
interest is a low-weight observable.

\section{Where \OBQF{} does \emph{not} transfer}
\label{sec:limits}

For completeness we also state the boundaries of the method, since this
is also part of an honest transferability discussion.
\begin{itemize}[leftmargin=1.2em,topsep=2pt,itemsep=2pt]
  \item \textbf{Continuous spectra without a gap.} Dense linear
        regression with an arbitrary covariance matrix, or a stiff PDE
        with no clear discretisation gap, will spread amplitude across
        the single phase bit and incur $\mathcal{O}(1)$ filter leakage.
        Multi-bit phase estimation or full QSVT is then required.
  \item \textbf{Tomography-bound answers.} If the answer is the full
        solution vector $\mathbf{x}$ rather than $\bra{\mathbf{x}}O\ket{\mathbf{x}}$
        for a low-weight $O$, the readout cost is
        $\mathcal{O}(N)$ regardless of the inversion step
        \cite{Aaronson:2015:readingthefineprint}, and any quantum speed
        up is erased -- this is precisely the issue Reviewer~2 raised
        with reference to \cite{Tang:2019} and we addressed in our
        Major-1 reply.
  \item \textbf{Expensive state preparation.} If $\ket{\mathbf{b}}$
        cannot be prepared without QRAM, condition (S2) fails and the
        speed-up is again lost.  Our tracking encoder happens to satisfy
        (S2) by Hadamards alone; transferred problems should be checked
        case by case.
  \item \textbf{Non-Hermitian or unbounded operators.} \OBQF{} requires
        Hamiltonian simulation $e^{-i\mathbf{A}t}$. Non-normal matrices
        (e.g.\ certain advection-dominated PDE discretisations) need
        block encoding via QSVT instead.
\end{itemize}

\section{Recommendation for the manuscript}
\label{sec:recommendation}

We propose adding to the manuscript a single short subsection,
``\emph{Generality of the spectral-filter mechanism}'', that:
(i)~states the three structural ingredients (S1)--(S3);
(ii)~cites the QSVT/IPE generalisation
\cite{Gilyen:2019,Martyn:2021,Kitaev:1995};
(iii)~lists by name two or three of the families F1--F6 above as
illustrative examples (associative memory, QUBO with structured
penalties, regularised kernel methods);
(iv)~explicitly states the limits of Section~\ref{sec:limits}.

This addresses Reviewer~2's ``too tailored'' concern without overclaiming
generality: the method is not a universal quantum solver, but the
structural conditions under which it works are common, well-defined,
and verifiable in advance by classical spectral analysis.

\section{Proposed reviewer reply}
\label{sec:reply}

\noindent
The techniques presented in this work were developed to make TrackHHL
executable on quantum hardware by reducing qubit count, two-qubit gate
count, and circuit depth -- the standard bottlenecks of HHL-based
pipelines. As the reviewer correctly notes, the optimizations exploit
the geometrical structure of the tracking problem and the numerical
features of our quadratic Hamiltonian, and this degree of
problem-specificity is typical of HHL applications, since HHL is not
currently regarded as a general-purpose replacement for classical
matrix inversion. We nevertheless believe that the methods transfer to
a well-defined family of problems, and we make the criterion for
transferability mathematical rather than application-driven.

Concretely, our one-bit quantum filter (\OBQF{}) solves
$\mathbf{A}\,\mathbf{x} = \mathbf{b}$ on the subspace where it is
informative, provided the input data $(\mathbf{A}, \mathbf{b}, O)$
satisfy the following structural conditions:
\begin{enumerate}[label=(\textbf{M\arabic*}),leftmargin=2.6em,topsep=2pt,itemsep=2pt]
  \item $\mathbf{A}\in\mathbb{R}^{N\times N}$ is symmetric positive
        semi-definite, $s$-sparse with $s = \mathcal{O}(\mathrm{polylog}\,N)$,
        and admits an efficient Hamiltonian-simulation oracle for
        $e^{-i\mathbf{A}t}$ (in our case, the sparse-segment encoder of
        Sec.~III).
  \item The spectrum of $\mathbf{A}$ decomposes as
        $\sigma(\mathbf{A}) = \sigma_\mathrm{sig} \cup \sigma_\mathrm{noise}$
        with
        $\sigma_\mathrm{noise} \subseteq [\lambda_n - \delta,
        \lambda_n + \delta]$ around an analytically known eigenvalue
        $\lambda_n$, and a finite spectral gap
        $\Delta = \mathrm{dist}(\sigma_\mathrm{sig},\sigma_\mathrm{noise}) > 0$
        with $\Delta \gg \delta$. The condition number on
        $\sigma_\mathrm{sig}$ is bounded by a problem-independent
        constant $\kappa_\mathrm{sig}=\mathcal{O}(1)$.
  \item $\ket{\mathbf{b}}$ is preparable by an
        $\mathcal{O}(\mathrm{polylog}\,N)$-depth circuit without QRAM
        (in our case $\ket{\mathbf{b}} = H^{\otimes n_s}\ket{0}$, the
        uniform superposition).
  \item The quantity of interest is
        $\bra{\mathbf{x}} O\ket{\mathbf{x}}/\|\mathbf{x}\|^{2}$ for an
        observable $O$ supported on $\mathcal{O}(1)$ qubits, e.g.\ a
        single-mode population $O = \ket{j}\bra{j}$ or a low-rank
        projector $O = \sum_{j\in S}\ket{j}\bra{j}$ with $|S|$ bounded
        independently of $N$.
\end{enumerate}
Whenever (M1)--(M4) hold, the same circuit applies up to a
problem-specific Hamiltonian-simulation block, with the same
$\mathcal{O}(\sqrt{N}\log N)$ gate budget reported in our paper.
This is the regime considered by Rebentrost et al.\
\cite{Rebentrost:2018:hopfield}, from which TrackHHL is inspired:
optimization of a Hopfield network is recast as a sparse linear problem
solved by HHL, with an exponential advantage provided that efficient
Hamiltonian simulation and readout subroutines are available, and that
the linear system can be suitably regularized. Our \OBQF{} construction
sharpens this picture: \emph{regularization} in their language is
exactly the noise-eigenvalue pinning of (M2), and \emph{efficient
readout} is exactly the low-weight observable of (M4). When both hold,
the multi-qubit phase register of standard HHL collapses to a single
clock qubit, as the only spectral information that has to be resolved
is the binary signal/noise label.

The structural conditions (M1)--(M4) define a family of problems
mathematically rather than by domain. Within this family one finds, in
addition to the RNA-reconstruction example outlined by Rebentrost et
al.: Hopfield-style associative memories and pattern denoising, where
$\mathbf{A}$ is a Hebbian coupling matrix whose spurious-pattern
spectrum concentrates on a known loading-fraction-dependent eigenvalue
\cite{Hopfield:1982,Amit:1985}; QUBO/Ising encodings of constrained
combinatorial problems whose conflict-penalty term sets $\lambda_n$ by
construction (max-cut, graph colouring, vertex cover, set cover, MIS;
see \cite{Lucas:2014}), where \OBQF{} returns a marginal of the
low-energy subspace and is complementary to QAOA \cite{Farhi:2014} and
quantum annealing \cite{Kadowaki:1998}; graph-Laplacian linear systems
$(\mathbf{L}+\alpha\mathbf{1})\mathbf{x}=\mathbf{b}$ on graphs with
$k$ well-separated communities, whose Laplacian spectrum has $k$ small
eigenvalues separated from a clearly defined bulk
\cite{vonLuxburg:2007,Apers:2022,Childs:2017}, with applications in
seeded community detection and effective-resistance queries; regularized
kernel systems $(\mathbf{K}+\lambda\mathbf{1})\boldsymbol{\alpha}
= \mathbf{y}$ in least-squares SVMs, ridge regression, and
Gaussian-process regression \cite{Williams:2006,Rebentrost:2014:qsvm},
where the regularizer pins $\lambda_n=\lambda$ and the prediction at a
test point is by construction a local observable on the output
register; and discretized linear PDEs
$\mathbf{L}_h\mathbf{u}=\mathbf{f}$ with a clear spectral gap between
global and discretization-scale modes \cite{Berry:2014:qode,Berry:2017},
provided only a local functional of $\mathbf{u}$ is read out -- a class
that includes track extrapolation through piecewise-uniform field maps.

Conversely, (M1)--(M4) are sharp, and the method does \emph{not}
transfer to problems with continuous, gapless spectra
($\Delta \to 0$, where filter leakage becomes $\mathcal{O}(1)$); to
tasks requiring full-state tomography (where the
$\mathcal{O}(N)$ readout cost erases any speed-up, as discussed in our
reply to Reviewer~2 Major-1
\cite{Aaronson:2015:readingthefineprint,Tang:2019}); to settings
where state preparation requires QRAM and (M3) fails; or to
non-normal $\mathbf{A}$, for which Hamiltonian simulation must be
replaced by block encoding via QSVT \cite{Gilyen:2019,Martyn:2021}.
We expect that, as in the present work, application to any of these
problems will still require adaptation to the specific structure and
constraints of the system at hand, but the conditions under which the
adaptation is worthwhile are well-defined and can be verified
classically -- by sparse spectral analysis of $\mathbf{A}$ -- before
any quantum circuit is written.

\bibliographystyle{unsrt}
\bibliography{refs}

\end{document}
